\chapter{Discussion}
\label{chap:discussion}

The results of this work show that publicly available GTFS and GTFS-Realtime data can be transformed into sustainability-relevant object-centric event logs. The created event logs contain attributes such as distance, GPS coordinates, delay, occupancy, speed, and information about the transport mode. These attributes provide the basis for sustainability and emission analysis.

However, the approach also has several limitations. One of the main challenges is the availability and quality of GTFS-Realtime data. While static GTFS data are available for many transport agencies, the usefulness of the approach strongly depends on the available real-time information. Some providers additionally offer attributes such as congestion level, which indicate whether a vehicle is operating under normal traffic conditions or in congested traffic. Such information could improve the accuracy of the emission estimation, but it was not available in the selected dataset. Furthermore, the available information differs between datasets. For example, other contributions, such as \cite{data10080119}, provide additional timestamps that allow a more detailed reconstruction of operational activities, including direction changes or shift-related activities. In our dataset, these activities mainly serve as markers instead of representing operational processes. Nevertheless, the selected dataset was well suited for this work because it provides vehicle IDs, delays, and occupancy information.

Another limitation concerns the information available about the vehicle fleet. Accurate emission estimation requires information about the vehicle model, fuel type, fuel consumption, and passenger capacity. Since these data are not part of the GTFS standard, external sources have to be used to estimate these values. Furthermore, occupancy is only available as status categories instead of the exact number of passengers. As a result, per-passenger emission calculations require additional assumptions.

Another limitation concerns the heuristic nature of the proposed event log construction. Several parts of the transformation pipeline rely on domain-specific assumptions to derive operational semantics from GTFS data, as this information is not explicitly available in the source data. Consequently, different heuristic choices may lead to different event logs and, therefore, different sustainability analysis results. In particular, the Operational State Heuristic (H5) depends on a temporal threshold that separates transitions within passenger service from transitions between operational phases. The chosen threshold directly influences the resulting distribution of activity types. A lower or higher threshold can lead to different numbers of \textit{parking}, \textit{begin\_shift}, and \textit{change\_direction} event types. This also affects the emission analysis, since different activity types define different segment boundaries and therefore influence the aggregation of distances, travel durations, and emissions.

Moreover, the semantic interpretation of the derived activity types influences the resulting emission estimates. In this work, the activity types \textit{begin\_shift} and \textit{parking} serve as markers for deadhead movements, for example between the depot and the first passenger stop or from the last passenger stop back to the depot, rather than explicitly representing these movements themselves. However, the selected GTFS dataset does not contain GPS trajectories for these deadhead trips. Consequently, their emissions cannot be estimated directly and can only be approximated or remain unaccounted for. Additional operational information, such as depot locations, scheduled driver breaks, vehicle rotations, or legally required layover times, could improve the identification of operational phases and enable a more reliable distinction between passenger trips, deadhead trips, parking periods, and layovers. Since these operational phases exhibit different driving behavior and emission characteristics, incorporating such domain knowledge would also improve the accuracy of the resulting emission estimates.

The evaluation also showed that both the distance-based and the duration-based approaches only provide approximations of the actual emissions. More accurate emission estimates require more detailed vehicle position data, providing speed information at a finer temporal or spatial granularity. However, these data were not available in the selected dataset. Since the proposed pipeline transforms the available GTFS and GTFS-Realtime data into an object-centric event log, information that is missing in the source data is also unavailable in the resulting OCEL. If such information becomes available in future GTFS-Realtime datasets, it could be integrated into the object-centric event log as additional event attributes. Instead of representing one segment between two consecutive stops with a single average speed, the segment could be divided into smaller sections with their own speed information. This would make it possible to distinguish between different driving situations, such as stop-and-go traffic, slow urban traffic, or free-flowing traffic. Based on these observations, different driving activities, for example \textit{drive-20km/h}, \textit{drive-30km/h}, or \textit{drive-50km/h}, could be derived to estimate emissions more accurately.

Although the approach was developed for public transportation, it could also be applied to other transport domains, such as logistics, where vehicles also follow predefined routes and stop at known locations. In contrast, it is less suitable for domains with highly dynamic and irregular movements, where no fixed routes or stops exist that can be used to reconstruct trips and estimate emissions.

Despite these limitations, this work demonstrates that GTFS and GTFS-Realtime data can be transformed into sustainability-relevant object-centric event logs. Since GTFS covers different transport modes, including buses, trains, trams, and ferries, the proposed pipeline can also serve as a basis for future sustainability analyses in public transportation.